% user_guide.tex — complete reference for every user-facing feature
\chapter{User Guide}
\label{chap:user-guide}

This chapter is the authoritative reference for every user-facing command,
configuration option, document format, and error condition exposed by
\cgspkg{}.

% -----------------------------------------------------------------------
\section{CLI Overview}

The single entry-point is the \cgscli{} command.

\begin{lstlisting}[language=bash]
$ cgitsync <command> [options]
\end{lstlisting}

Running \cgscli{} without a command prints a help summary and exits with
code~0.

% -----------------------------------------------------------------------
\section{Commands}

\subsection{\texttt{validate}}

\begin{lstlisting}[language=bash]
$ cgitsync validate <spec.cgs>
\end{lstlisting}

Parses the \texttt{.cgs} file, builds the dependency-tree registry, and
prints a one-line summary for each repository node. Exits~0 on success,
non-zero on parse or validation error.

\textbf{Output fields:} \texttt{name}, \texttt{lifecycle\_state},
\texttt{complete}, \texttt{node\_count}.

\subsection{\texttt{tree}}

\begin{lstlisting}[language=bash]
$ cgitsync tree <spec.cgs>
\end{lstlisting}

Renders the dependency graph of the spec as a text tree. Useful for a quick
visual sanity-check of parent/child relationships.

\subsection{\texttt{describe}}

\begin{lstlisting}[language=bash]
$ cgitsync describe <snapshot.gts>
\end{lstlisting}

Reads a \texttt{.gts} snapshot and prints its contents as pretty-printed
JSON. Useful for inspecting saved tree state after a \texttt{clone} or
\texttt{freeze\_release} operation.

\subsection{Planned Commands (not yet implemented)}

The following commands are defined in the public API contract but are not yet
implemented. Each returns exit code~2 with the message
\texttt{"not implemented yet"}.

\begin{itemize}
  \item \texttt{clone} — bootstrap a full repository tree from a \texttt{.cgs} spec.
  \item \texttt{checkout} — switch the whole tree to a named ref.
  \item \texttt{commit} — stage and commit changes across the tree.
  \item \texttt{push} — push all repos to their remotes.
  \item \texttt{tag} — create a named tag across the tree.
  \item \texttt{freeze\_release} — write a \texttt{.gts} snapshot and tag the release.
  \item \texttt{restart} — re-synchronise a tree from an existing \texttt{.gts} snapshot.
  \item \texttt{launch\_release} — check out a previously frozen release.
\end{itemize}

% -----------------------------------------------------------------------
\section{Document Formats}

\subsection{Local Authoring Spec (\texttt{.cgs})}

A TOML file that describes the repository tree to manage. It is the only
hand-authored input; it is never modified at runtime.

\textbf{Top-level sections:}

\begin{itemize}
  \item \texttt{[document]} — mandatory; contains \texttt{format\_version}.
  \item \texttt{[project]} — mandatory; project-level metadata (see
        Table~\ref{tab:project-fields}).
  \item \texttt{[[repos]]} — one entry per managed repository
        (see Table~\ref{tab:repo-fields}).
\end{itemize}

\begin{table}[h]
\centering
\begin{tabular}{lll}
\hline
\textbf{Field} & \textbf{Type} & \textbf{Description} \\
\hline
\texttt{name}           & string  & Human-readable project label \\
\texttt{repo\_name}     & string  & Git repository slug \\
\texttt{default\_branch}& string  & Default branch (e.g.\ \texttt{main}) \\
\texttt{gitprovider}    & string  & Provider (\texttt{github}, \texttt{gitlab}, \texttt{custom}) \\
\texttt{gitprovider\_url}& string & Required for \texttt{custom} provider \\
\hline
\end{tabular}
\caption{Project-level fields in \texttt{.cgs}}
\label{tab:project-fields}
\end{table}

\begin{table}[h]
\centering
\begin{tabular}{lll}
\hline
\textbf{Field} & \textbf{Type} & \textbf{Description} \\
\hline
\texttt{gitprovider}         & string  & Provider override for this repo \\
\texttt{project\_owner\_name}& string  & GitHub owner / GitLab group \\
\texttt{project\_name}       & string  & Repository name on the provider \\
\texttt{relative\_path}      & string  & Path relative to project root \\
\texttt{gitprovider\_url}    & string  & Required for \texttt{custom} provider \\
\hline
\end{tabular}
\caption{Per-repository fields in \texttt{.cgs}}
\label{tab:repo-fields}
\end{table}

\subsection{Git Tree State Snapshot (\texttt{.gts})}

A TOML file generated automatically by bootstrapping and release operations.
It must never be hand-edited. It captures the exact commit SHAs and lifecycle
state of the tree at a point in time.

\textbf{Top-level sections:}
\begin{itemize}
  \item \texttt{[document]} — format version, \texttt{generated\_at},
        \texttt{command\_origin}.
  \item \texttt{[project]} — project name and \texttt{root\_absolute\_path}.
  \item \texttt{[tree\_state]} — \texttt{lifecycle\_state}, \texttt{is\_ready},
        \texttt{registry\_complete}.
  \item \texttt{[[repo\_state]]} — one entry per repo; includes
        \texttt{commit\_sha}, \texttt{sync\_state}, and ref information.
\end{itemize}

\subsection{Planning / GOC Documents (\texttt{.goc})}

Used for project identity and planning metadata. Supported in TOML, YAML, or
JSON. YAML support requires the optional \texttt{PyYAML} dependency.

% -----------------------------------------------------------------------
\section{Python API}

\cgspkg{} can be used as a library. The public API is exported from the
top-level package:

\begin{lstlisting}[language=Python]
from ComplexGitSync import (
    ComplexGitSyncClient,
    CgsDocument,
    GtsDocument,
    GocDocument,
    GitProvider,
    AccessProtocol,
    RepoAddress,
    GitRepo,
    GitTree,
    Orchestre,
    DependencyTreeRegistry,
)
\end{lstlisting}

\subsection{\texttt{ComplexGitSyncClient}}

High-level client that mirrors the CLI commands. Instantiate with a parsed
\texttt{CgsDocument} and call the equivalent method:

\begin{lstlisting}[language=Python]
from ComplexGitSync import ComplexGitSyncClient, CgsDocument

doc = CgsDocument.from_toml(pathlib.Path("project.cgs").read_bytes())
client = ComplexGitSyncClient(doc)
summary = client.validate()
\end{lstlisting}

\subsection{Document Classes}

Each document class exposes six format factories:

\begin{lstlisting}[language=Python]
doc = CgsDocument.from_toml(raw_bytes)
doc = CgsDocument.from_json(raw_str)
doc = CgsDocument.from_yaml(raw_str)   # requires PyYAML

serialised = doc.to_toml()
serialised = doc.to_json()
serialised = doc.to_yaml()             # requires PyYAML
\end{lstlisting}

\subsection{\texttt{RepoAddress}}

Composes and parses remote repository URLs:

\begin{lstlisting}[language=Python]
from ComplexGitSync import RepoAddress, GitProvider, AccessProtocol

addr = RepoAddress(
    gitprovider=GitProvider.GITHUB,
    owner_name="my-org",
    project_name="my-repo",
)
print(addr.to_ssh())    # git@github.com:my-org/my-repo.git
print(addr.to_https())  # https://github.com/my-org/my-repo.git
\end{lstlisting}

% -----------------------------------------------------------------------
\section{Configuration and Environment}

\cgspkg{} reads no environment variables and uses no global configuration
file. All configuration is supplied through the \texttt{.cgs} file.

% -----------------------------------------------------------------------
\section{Error Reference}

\begin{description}
  \item[\texttt{ValueError}] Raised when a required field is missing or a
    field value is invalid (e.g.\ unknown \texttt{gitprovider},
    missing \texttt{gitprovider\_url} for a \texttt{custom} provider).
  \item[\texttt{NotImplementedError}] Raised when a planned but not yet
    implemented command is invoked via the Python API.
  \item[Exit code 1] General CLI failure (parse error, validation error).
  \item[Exit code 2] Command exists but is not yet implemented.
\end{description}

% -----------------------------------------------------------------------
\section{Troubleshooting}

\begin{description}
  \item[``\texttt{gitprovider\_url} is required for custom provider'']
    Add a \texttt{gitprovider\_url} field to the offending
    \texttt{[[repos]]} entry in your \texttt{.cgs} file.
  \item[\texttt{validate} exits non-zero despite a correct-looking file]
    Check that \texttt{[document]} contains \texttt{format\_version = "1.0"}
    and that every \texttt{[[repos]]} block includes
    \texttt{project\_owner\_name} and \texttt{project\_name}.
  \item[PyYAML import error when using \texttt{.goc} files]
    Install the optional extra: \texttt{uv pip install "ComplexGitSync[yaml]"}.
\end{description}
